% Theorem 1 for error rate \theta, as \theta decreases the volume of Rashomon set - the total number of models that achieve error \theta - decreases monotonically (maybe with somp probablity if we cannot prove it for certainty)

% Theorem 2 for error rate \theta, P and Q are two models randomly drawn from the rashomon set, as \theta dcreases, KL(P||Q) decrease (probablistically, or taking the expectation of KL) 


% --------------------------------------------------------------
% This is all preamble stuff that you don't have to worry about.
% Head down to where it says "Start here"
% --------------------------------------------------------------
 
\documentclass[12pt]{article}
 
\usepackage[margin=1in]{geometry} 
\usepackage{amsmath,amsthm,amssymb,dsfont,physics,scrextend}
\usepackage{fancyhdr}
\usepackage{xcolor}
\pagestyle{fancy}

 
\newcommand{\N}{\mathbb{N}}
\newcommand{\Z}{\mathbb{Z}}
\newcommand{\I}{\mathbb{I}}
\newcommand{\R}{\mathbb{R}}
\newcommand{\Q}{\mathbb{Q}}
\renewcommand{\qed}{\hfill$\blacksquare$}
\let\newproof\proof
\renewenvironment{proof}{\begin{addmargin}[1em]{0em}\begin{newproof}}{\end{newproof}\end{addmargin}\qed}
% \newcommand{\expl}[1]{\text{\hfill[#1]}$}
 
\newenvironment{theorem}[2][Theorem]{\begin{trivlist}
\item[\hskip \labelsep {\bfseries #1}\hskip \labelsep {\bfseries #2.}]}{\end{trivlist}}
\newenvironment{lemma}[2][Lemma]{\begin{trivlist}
\item[\hskip \labelsep {\bfseries #1}\hskip \labelsep {\bfseries #2.}]}{\end{trivlist}}
\newenvironment{problem}[2][Problem]{\begin{trivlist}
\item[\hskip \labelsep {\bfseries #1}\hskip \labelsep {\bfseries #2.}]}{\end{trivlist}}
\newenvironment{exercise}[2][Exercise]{\begin{trivlist}
\item[\hskip \labelsep {\bfseries #1}\hskip \labelsep {\bfseries #2.}]}{\end{trivlist}}
\newenvironment{reflection}[2][Reflection]{\begin{trivlist}
\item[\hskip \labelsep {\bfseries #1}\hskip \labelsep {\bfseries #2.}]}{\end{trivlist}}
\newenvironment{proposition}[2][Proposition]{\begin{trivlist}
\item[\hskip \labelsep {\bfseries #1}\hskip \labelsep {\bfseries #2.}]}{\end{trivlist}}
\newenvironment{corollary}[2][Corollary]{\begin{trivlist}
\item[\hskip \labelsep {\bfseries #1}\hskip \labelsep {\bfseries #2.}]}{\end{trivlist}}

\newtheorem{definition}{Definition}

\begin{document}
 
% --------------------------------------------------------------
%                         Start here
% --------------------------------------------------------------

\lhead{Misuse Theorems}
\chead{Ronilo Ragodos}
\rhead{\today}
 
% \maketitle
% SHOULD LOOK AT RASHOMON LOSS DISTRIBUTION

\section{Notation and Definitions}
 
  

Suppose we are given:
\begin{itemize}
    \item a multi-class classification dataset $\mathcal{D}=(X,Y)\subseteq \mathcal{X}\times \mathcal{Y}$ with and $n$ instances and $c$ categories
    \item hypothesis space $\mathcal{H}$ of probabilistic classifiers of $\mathcal{D}$, each parameterized by some $\phi \in \Phi \subset \mathbb{R}^d$ and approximating the true conditional distribution $P_{Y\mid X=\mathbf{x}_i}$. We write $\mathcal{H}=\{h_\phi :\mathcal{X} \rightarrow \Delta_c\}$ where $\Delta_c = \{(r_1,\ldots,r_c)\in [0,1]^c\mid \sum_{i=1}^c r_i = 1\}$ is the $c$-dimensional probability simplex.
    % \item cross entropy loss $l:\Delta_c \times \Delta_c \rightarrow \mathbb{R}^+$ used to evaluate the hypotheses
    % \item population risk $L(h_\theta) = \mathds{E}_{P_{Y\mid X}}[l(h_\phi(X),Y)]$
    % \item empirical risk $\hat{L}(h_\phi) = \frac{1}{n} \sum_{i=1}^n l(h_\phi(X),Y))$
\end{itemize}

% \begin{definition}[Rashomon Ratio]
%     Define $\mathcal{H}_\epsilon =\{h_\phi \in \mathcal{H} \mid L(h_\phi) \leq \epsilon\}$ and let $\lambda$ be a probability measure on $\mathbb{R}^d$. Then the \emph{Rashomon ratio} (with respect to $\lambda$) is given by $\frac{\lambda(\mathcal{H}_\epsilon)}{\lambda(\mathcal{H})}$. 
% \end{definition}

\pagebreak

\section{Theorems Relating to Decision Process}



Let $P,Q \in \mathcal{H}_{\theta_n}$ be two probabilistic classifiers. Let $(X,Y)$ be a random variable representing samples of size $n$ from $\mathcal{X}\times \mathcal{Y}$. Let $\hat{Y}$ be a random vector of outputs of $P$, so that $\hat{Y}_n \sim \text{Bernoulli}(P(\mathbf{x}_n), 1-P(\mathbf{x}_n))$.


Suppose that $\mathds{P}[Y = \hat{Y}] = .95$

\begin{proposition}{A}
    Suppose that $\mathcal{H}$ consists of logistic regressions and $\mathcal{G}$ itself is a logistic regression. Suppose further that $\mathcal{X}$ has $D$ real valued independent normally distributed features. 

    Claim: the following terms are non increasing as $\theta$ is decreased.
    \begin{equation}
        \mathds{E}_{P\sim \mathcal{H}_{\theta}}[\norm{\hat{\beta} - \beta^*}]
    \end{equation}
\end{proposition}

The coefficient errors of logistic regression are given by the diagonal terms of the covariance matrix. This is a wishart matrix. Therefore the coefficient errors are chi-square random variables. Hence $\hat{\beta} - \beta^*$ is a chi-square random vector. Chi-square random variables are sub-exponential so we may use the bound:
% https://arxiv.org/pdf/1804.00205
\begin{equation*}
    \mathds{E}\left[\norm{\hat{\beta}-\beta^*}_2\right] \leq c \max_d \mid \beta_d - \beta^*_d \mid
\end{equation*}


flow of ideas for next step
\begin{itemize}
    %\item $\mathds{P}[Y=\hat{Y}]$ being high means usually $\mid P(x)-G(x)\mid$ is small (NOT TRUE)
    \item Suppose total variation distance is small; $\delta(P,G)=\sup_{E}\mid P(E) - G(E)\mid = \epsilon_1$
    \item Sigmoid has bounded derivative so it is lipschitz
    \item \textcolor{red}{SUPPOSE UNIFORM LIP implies coordinate lipschitz}
    \item if sigmoid is coordinate wise lipschitz in the coefficients (for fixed x), then $|\hat{\beta}_d - \beta^*_d| \leq L_d |P(x)-G(x)| \leq L_d \epsilon_1$
    \item So if $L_d$ is small, get result???
    \item Also, then $
    \mathds{E}\left[\norm{\hat{\beta}-\beta^*}_2\right] \leq c \max_d L_d |P(x)-G(x)| \leq c \max_d L_d \epsilon_1$
\end{itemize}

can also do some tail bounds?:
\begin{gather*}
    \mathds{P}\left[\norm{\hat{\beta}-\beta^*}_2 >  \mathds{E}\left[\norm{\hat{\beta}-\beta^*}_2\right]\right] \\
    \sim \mathds{P}\left[\norm{\hat{\beta}-\beta^*}_2 >  c \max_d \mid \beta_d - \beta^*_d\right] \\
    \mathds{P}\left[\norm{\hat{\beta}-\beta^*}_2 >  \epsilon\right] \\
\end{gather*}
\pagebreak

\section{Theorems Relating to Output Only}


  \begin{lemma}{1}
    For any measure $\lambda$ on $\mathbb{R}^d$, if we define
    \begin{equation*}
        \mathcal{H}_\theta = \{h_\phi \in \mathcal{H} \mid \theta - \delta \leq L(h_\phi)\leq \theta\}
    \end{equation*}
    where $\delta \leq \theta$ is fixed, then $\lambda(\mathcal{H}_\theta)$ is non-increasing as $\theta\rightarrow \delta$.
 \end{lemma}

\begin{proof}
    Measures are monotone, so it suffices to show that for a pair $\theta' < \theta$, 
    \begin{equation*}
        \mathcal{H}_{\theta'} \subseteq \mathcal{H}_\theta
    \end{equation*}

    By way of contradiction, suppose that there exists a pair $\theta' < \theta$ such that  $\mathcal{H}_{\theta'} \setminus \mathcal{H}_\theta \neq \emptyset$. Let $h^*$ be a member of $\mathcal{H}_{\theta'} \setminus \mathcal{H}_\theta$. $h^*$ must be such that
    \begin{gather*}
        L(h^*) \in [\theta'-\delta,\theta'] \text{ and } L(h^*) \notin [\theta-\delta,\theta]
    \end{gather*}
    But since $\theta > \theta'$, $[\theta-\delta,\theta]\supset [\theta'-\delta,\theta']$ and so we reach a contradiction.

    Therefore, it must be the case that for any $\theta' < \theta$, $\mathcal{H}_{\theta'} \subseteq \mathcal{H}_\theta$. It follows by the monotonicity of measure that $\lambda(\mathcal{H}_{\theta'})\leq \lambda(\mathcal{H}_\theta)$. 
\end{proof}

\begin{theorem}{1}
        Suppose that we have a binary classification dataset $\mathcal{D}\subset \mathcal{X}\times [0,1]$ with $D$ features and $N$ instances whose labels were sampled according to a ground truth distribution $G:X \rightarrow [0,1]$. 

    We define for $\theta \geq H(G)$ the hypothesis classes $\mathcal{H}_\theta$ of probabilistic classifiers of $D$ with empirical loss upper bounded by $\theta$:
    \begin{equation*}
        \mathcal{H}_\theta = \{h \in \mathcal{H} \mid L(h) \leq \theta\}
    \end{equation*}

    Assume that the distribution $G$ has the property that for any $\theta$ and any $P \in \mathcal{H}_\theta$, $L(G) \leq L(P)$. Further assume that $G$ is unique among hypotheses in having this property. 

    Let $(\theta_n)$ be an decreasing sequence of non-negative numbers bounded below by $H(G)$. Then, if for any pair $n,n+1$ we have $\mathds{E}_{P\sim \mathcal{H}_{\theta_{n+1}}}[H(P)] - \mathds{E}_{P'\sim \mathcal{H}_{\theta_{n}}}[H(P')] \geq 0$, the following sequence of terms is non-increasing in $n$. Here we assume the expectations are computed with respect to some probability measure on $\mathcal{H}$.
    \begin{equation*}
        \mathds{E}_{P\sim \mathcal{H}_{\theta_{n+1}}}[D_{KL}(P || G)] - \mathds{E}_{P\sim \mathcal{H}_{\theta_n}}[D_{KL}(P || G)] 
    \end{equation*}
\end{theorem}

\begin{proof}
    \begin{gather*}
         \mathds{E}_{P\sim \mathcal{H}_{\theta_{n+1}}}[D_{KL}(P || G)] - \mathds{E}_{P'\sim \mathcal{H}_{\theta_n}}[D_{KL}(P || G)] =  \mathds{E}_{P\sim \mathcal{H}_{\theta_{n+1}}}[L(P) - H(P)] - \mathds{E}_{P'\sim \mathcal{H}_{\theta_n}}[L(P')-H(P')]  \\
    \end{gather*}
    Thus we seek to satisfy
    \begin{gather*}
        \mathds{E}[L(P)]-\mathds{E}[L(P')] \leq \mathds{E}[H(P)]-\mathds{E}[H(P')]
    \end{gather*}
    We know that the LHS is $\theta_{n+1}-\theta_n < 0$, so since we hypothesized that $\mathds{E}[H(P)]-\mathds{E}[H(P')] \geq 0$ the inequality is assured. 
\end{proof}

\begin{proposition}{1}
    Suppose that we have a binary classification dataset $\mathcal{D}\subset \mathcal{X}\times [0,1]$ with $D$ features and $N$ instances whose labels were sampled according to a ground truth distribution $G:X \rightarrow [0,1]$. 

    We define for $\theta \geq H(G)$ the hypothesis classes $\mathcal{H}_\theta$ of probabilistic classifiers of $D$ with empirical loss upper bounded by $\theta$:
    \begin{equation*}
        \mathcal{H}_\theta = \{h \in \mathcal{H} \mid \hat{L}(h;D) \leq \theta\}
    \end{equation*}

    Assume that the distribution $G$ has the property that for any $\theta$ and any $P \in \mathcal{H}_\theta$, $\hat{L}(G;D) \leq \hat{L}(P;D)$. Further assume that $G$ is unique among hypotheses in having this property. 

    Let $(\theta_n)$ be an decreasing sequence of non-negative numbers bounded below by $H(G)$ such that for every $n$, we may construct a finite sample $\hat{\mathcal{H}}_{\theta_n}$ that is non-empty. Define $M_n = |\hat{\mathcal{H}}_{\theta_n}|$. Then there are constants $B_{n+1},B_{n\setminus n+1}>0$ such that the following bound holds.

    \begin{equation*}
        \mathds{E}_{P\sim \hat{\mathcal{H}}_{\theta_{n+1}}}[D_{KL}(P || G)] - \mathds{E}_{P\sim \hat{\mathcal{H}}_{\theta_n}}[D_{KL}(P || G)] \leq M_{n+1}B_{n+1}(\frac{1}{M_{n+1}} - \frac{1}{M_n}) - \frac{1}{M_n}(M_n - M_{n+1} + 1)B_{n\setminus n+1}
    \end{equation*}
\end{proposition}

\begin{proof}
    % By Lemma 1, $(\hat{\mathcal{H}}_{\theta_n})$ is a non-increasing sequence of sets. Therefore, it is the case that:
    % \begin{equation*}
    %     \lim_{n\rightarrow \infty} \hat{\mathcal{H}}_{\theta_n} = \bigcap_{n\geq 1} \hat{\mathcal{H}}_{\theta_n}
    % \end{equation*}
    % By definition of $G$, $G \in \bigcap_{n\geq 1} \hat{\mathcal{H}}_{\theta_n}$. Further, because of its uniqueness $\bigcap_{n\geq 1} \hat{\mathcal{H}}_{\theta_n} = \{G\}$ (this can be proved by reductio). 
    % It then follows that
    % \begin{equation*}
    %     \lim_{n\rightarrow \infty} \mathds{E}_{P \sim \hat{\mathcal{H}}_{\theta_n}}[D_{KL}(P||G)] = D_{KL}(G||G) = \hat{L}(G) - \hat{H}(G)
    % \end{equation*}
    Let us begin by considering an arbitrary pair $\theta_{n+1}$ and $\theta_n$ for some $n$. We need to give an upper bound to the following expression
    \begin{gather*}
        \mathds{E}_{P\sim \hat{\mathcal{H}}_{\theta_{n+1}}}[D_{KL}(P || G)] - \mathds{E}_{P\sim \hat{\mathcal{H}}_{\theta_n}}[D_{KL}(P || G)] 
    \end{gather*}
    We are assuming that the expectations are computed with respect to uniform distributions.  So, we have
    \begin{gather*}
       \frac{1}{M_{n+1}} \sum_{i=1}^{M_{n+1}} D_{KL}(P_i||G) - \frac{1}{M_n}\sum_{j=1}^{M_n} D_{KL}(P_j||G)
    \end{gather*}
    We may write this difference of sums in terms of the summands held in common and those that are exclusive to the larger set $\hat{\mathcal{H}}_{\theta_n}$:
    \begin{gather*}
        \sum_{i=1}^{M_{n+1}} D_{KL}(P_i||G)(\frac{1}{M_{n+1}}-\frac{1}{M_n}) - \frac{1}{M_n}\sum_{j=M_{n+1}+1}^{M_n} D_{KL}(P_j||G)
    \end{gather*}
    Let $B_{n+1}$ be an upper bound of $D_{KL}(\cdot||G)$ on $\hat{\mathcal{H}}_{\theta_{n+1}}$ and $B_{n\setminus n+1}$ be a lower bound of $D_{KL}(\cdot||G)$ on $\hat{\mathcal{H}}_{\theta_n} \setminus \hat{\mathcal{H}}_{\theta_{n+1}}$. This yields the following inequality:
        \begin{gather*}
        \sum_{i=1}^{M_{n+1}} D_{KL}(P_i||G)(\frac{1}{M_{n+1}}-\frac{1}{M_n}) - \frac{1}{M_n}\sum_{j=M_{n+1}+1}^{M_n} D_{KL}(P_j||G) \\
        \leq \sum_{i=1}^{M_{n+1}} B_{n+1}(\frac{1}{M_{n+1}}-\frac{1}{M_n}) - \frac{1}{M_n}\sum_{j=M_{n+1}+1}^{M_n} B_{n\setminus n+1}\\
        = M_{n+1}B_{n+1}(\frac{1}{M_{n+1}} - \frac{1}{M_n}) - \frac{1}{M_n}(M_n - M_{n+1} + 1)B_{n\setminus n+1}
    \end{gather*} 
    as required.
\end{proof}

\begin{corollary}{1}
    Define 
    \begin{gather*}
        P_{n+1}:=\arg \min_{P\in \hat{\mathcal{H}}_{\theta_{n+1}\setminus \{G\}}} D_{KL}(P||G) \\
        P_{n\setminus n+1} := \arg \max_{P\in \hat{\mathcal{H}}_{\theta_{n}\setminus \hat{\mathcal{H}}_{\theta_{n+1}}}} D_{KL}(P||G) \\
    \end{gather*}
    
    If, in the settings of Proposition 1 we additionally assume that the sequence $\hat{\mathcal{H}}_{\theta_n}$ can be constructed so that for any $n,n+1$, the following hold:
    % "loss reduction exceeds entropy reduction"
    \begin{itemize}
        \item $\mid \hat{\mathcal{H}}_{\theta_n} \mid > \mid \hat{\mathcal{H}}_{\theta_{n+1}}\mid$
        \item $L(P_{n\setminus n+1}) - L(P_{n+1}) \geq H(P_{n\setminus n+1}) - H(P_{n+1})$
    \end{itemize}
    Then the sequence indexed by $n$
    \begin{equation*}
        (\mathds{E}_{P\sim \hat{\mathcal{H}}_{\theta_n}}[D_{KL}(P || G)])
    \end{equation*}
    is a non-increasing sequence. %furthermore it converges to H(G)

    \bigskip

    Note that since we should often be able to safely assume that $L(P_{n\setminus n+1}) - L(P_{n+1}) \geq 0$, in such cases a more relevant hypothesis is that $H(P_{n\setminus n+1}) - H(P_{n+1}) \leq 0$, i.e. that the distribution in $\mathcal{H}_{\theta_n}\setminus\mathcal{H}_{\theta_{n+1}}$ that is furthest from $G$ in terms of KL divergence has less entropy than the closest distribution to $G$ in $\mathcal{H}_{\theta_{n+1}}$ does. 
\end{corollary}

\begin{proof}
    The previous proposition stated the following bound:
    \begin{equation*}
        \mathds{E}_{P\sim \hat{\mathcal{H}}_{\theta_{n+1}}}[D_{KL}(P || G)] - \mathds{E}_{P\sim \hat{\mathcal{H}}_{\theta_n}}[D_{KL}(P || G)] \leq M_{n+1}B_{n+1}(\frac{1}{M_{n+1}} - \frac{1}{M_n}) - \frac{1}{M_n}(M_n - M_{n+1} + 1)B_{n\setminus n+1}
    \end{equation*}
    The above is negative whenever $B_{n\setminus n+1} \geq B_{n+1} > 0$ and $M_n > M_{n+1}$. The latter constraint was assumed. 
    Tight choices of $B_{n+1}$ and $B_{n\setminus n+1}$ that satisfy the first constraint are:
    \begin{gather*}
        B_{n+1} = \min_{P \in \hat{\mathcal{H}}_{\theta_{n+1}} \setminus \{G\}} D_{KL}(P||G) \\
        B_{n\setminus n+1} = \max_{P \in \hat{\mathcal{H}}_{\theta_n} \setminus \hat{\mathcal{H}}_{\theta_{n+1}}} D_{KL}(P||G)
    \end{gather*}

    Our requirement then becomes $L(P_{n\setminus n+1}) - H(P_{n\setminus n+1}) \geq L(P_{n \setminus n+1}) - H(P_{n+1})$, which is true by hypothesis. 
    % TRY BOUNDING IN TERMS OF THE THETAS!!!!
\end{proof}

\begin{proposition}{2}
    $L(P_{n\setminus n+1}) - L(P_{n+1}) \geq H(P_{n\setminus n+1}) - H(P_{n+1})$ holds with probability...

    \item $L(P_{n\setminus n+1}) - L(P_{n+1}) \geq 0$ holds with probability ....

    \item \item $H(P_{n\setminus n+1}) - H(P_{n+1}) \leq 0$ holds with probability....
\end{proposition}



\begin{proposition}{TEMP}
    Suppose that we have a binary classification dataset $\mathcal{D}\subset \mathcal{X}\times [0,1]$ with $D$ features and $N$ instances whose labels were sampled according to a ground truth distribution $G:X \rightarrow [0,1]$. 
    
    For $\theta \geq 0$ if we define $\hat{\mathcal{H}}_\theta = \{h\in \mathcal{H} \mid \hat{L} \leq \theta\}$ to be the class of probabilistic classifiers of $\mathcal{D}$ with empirical loss at most $\theta$. Then we have the following three upper bounds:
    \begin{gather*}
        \mathds{E}_{P\sim \hat{\mathcal{H}}_\theta}[\widehat{D(P||G)}] \leq \mathds{E}_{P\sim \hat{H}_\theta}[\hat{L}(P)] - H(G) \\
        \mathds{E}_{P\sim \hat{\mathcal{H}}_\theta}[\widehat{D(P||G)}] \leq \hat{L}(\mathds{E}_{P\sim \hat{H}_\theta}[P]) - H(G) \\
        \mathds{E}_{P\sim \hat{\mathcal{H}}_\theta}[\widehat{D(P||G)}] \leq \theta - H(G) \\
    \end{gather*}
    %Furthermore, both upper bounds are non-increasing as $\theta \rightarrow 0$. 
\end{proposition}

\begin{proof}
We begin by decomposing KL divergence using its relation with cross entropy. 
\begin{gather*}
    \mathds{E}_{P\sim \hat{\mathcal{H}}_\theta}[\widehat{D(P||G)}] = \mathds{E}_{P\sim 
    \hat{\mathcal{H}}_\theta}[\hat{L}(P) - \hat{H}(P)] \\
    = \mathds{E}_{P\sim \hat{\mathcal{H}}_\theta}[\hat{L}(P)] - \mathds{E}_{P\sim \hat{\mathcal{H}}_\theta}[\hat{H}(P)] \\
\end{gather*}

%https://math.stackexchange.com/questions/1114589/entropy-is-upper-bounded-by-cardinality-of-the-random-variable
% Now may obtain the first upper bound in the theorem statement by upper bounding the second summand using Jensen's inequality and the fact that a uniform distribution maximizes entropy:
% \begin{gather*}
%     %\mathds{E}_{P\sim \mathcal{H}_\theta}[H(P)] \leq \log_2 |\mathcal{H}_\theta|
%     \hat{H}(P) = \mathds{E}_{x\sim X}[-\log(P(x))] \leq -\log(\mathds{E}[P(x)]) \\
%     = \log\left( \frac{1}{\mathds{E}[P(x)]} \right) \\ 
%     \leq \log |X| = \log N
% \end{gather*}
Since for random variables $A,B$ it is the case that $A \leq B \implies \mathds{E}[A] \leq \mathds{E}[B]$, we now have the first upper bound:
\begin{equation*}
     \mathds{E}_{P\sim \hat{\mathcal{H}}_\theta}[\widehat{D(P||G)}] \leq \mathds{E}_{P\sim \hat{\mathcal{H}}_\theta}[\hat{L}(P)] - \log N
\end{equation*}
Then the second upper bound in the theorem statement follows by applying Jensen's inequality to the above, noting that cross entropy is convex. 

The third inequality follows from the first and the definition of $\hat{\mathcal{H}}_\theta$. 

 \end{proof}

\begin{corollary}{TEMP}
    If the assumptions of proposition are modified so that $\mathcal{H}_\theta = \{h \in \mathcal{H}\mid \hat{L}(h) = \theta\}$ and we define a decreasing sequence $(\theta_n)$ such that for each $n$ there exists at least one $h \in \mathcal{H}$ such that $\hat{L}(h)=\theta_n$, and the sequence $(\mathds{E}_{P\sim \hat{\mathcal{H}}_{\theta_n}}[\hat{H}(P)])$ (indexed by $n$) is decreasing, then
    \begin{gather*}
        \mathds{E}_{P\sim \hat{\mathcal{H}}_{\theta_n}}[\widehat{D(P||G)}] = \theta_n - \mathds{E}_{P\sim \hat{\mathcal{H}}_{\theta_n}}[\hat{H}(P)]
    \end{gather*}
    and the sequence of these terms (indexed by $n$) is decreasing. In other words, the expected KL divergence between distributions in $\mathcal{H}_{\theta_n}$ and $G$ decreases with $n$. 
\end{corollary}
% 

\begin{proposition}{Hellinger Bound}
    %Given a binary classification dataset $\mathcal{D}=\mathcal{X}\times \mathcal{Y}$ with $D$ features and $N$ instances and an expected prediction error rate $\epsilon$, if we define $\mathcal{R}$ to be the class of probabilistic classifiers of $\mathcal{D}$ that achieve expected error rate $\epsilon$, then 
    Let $\lambda$ be \emph{any} probability measure on $\mathcal{H}$ and let $\delta,\theta$, $\mathcal{H}_\theta$ be as defined previously, and $D_H$ be the Hellinger distance. Then the terms
    \begin{equation*}
        \mathds{E}_{P,Q\sim \mathcal{H}_\theta}[D_H(P,Q)]
    \end{equation*}
    may be upper bounded such that the upper bound is non-increasing as $\theta \rightarrow \delta$. 
\end{proposition}

\begin{proof}
    Since Hellinger distance is a metric, we may invoke the triangle inequality and use Hellinger distance's relation to KL Divergence to obtain the following bounds. 
    \begin{gather*}
        D_H(P,Q) \leq D_H(P,G) + D_H(G,Q) \\
        = D_H(P,G) + D_H(Q,G) \\
        \leq \sqrt{D_{KL}(P||G)/2} + \sqrt{D_{KL}(Q||G)/2} \\
    \end{gather*}
    Taking expectations over $\mathcal{H}_{\theta}$, we have
    \begin{equation*}
        \mathds{E}_{P,Q\sim \mathcal{H}_\theta}[D_H(P,Q)] \leq \mathds{E}_{P\sim \mathcal{H}_\theta}[\sqrt{D_{KL}(P||G)/2}] + \mathds{E}_{Q\sim \mathcal{H}_\theta}[\sqrt{D_{KL}(Q||G)/2}] \\
        = 2\mathds{E}_{P\sim \mathcal{H}_\theta}[\sqrt{D_{KL}(P||G)/2}]
    \end{equation*}
    Since the square root function is monotonic, under the assumptions of Theorem 1 the RHS is non-increasing. 
\end{proof}

\begin{theorem}{Main}
    %Given a binary classification dataset $\mathcal{D}=\mathcal{X}\times \mathcal{Y}$ with $D$ features and $N$ instances and an expected prediction error rate $\epsilon$, if we define $\mathcal{R}$ to be the class of probabilistic classifiers of $\mathcal{D}$ that achieve expected error rate $\epsilon$, then 
    Let $\lambda$ be \emph{any} probability measure on $\mathcal{H}$ and let $\delta,\theta$, $\mathcal{H}_\theta$ be as defined previously, and $D$ be a divergence measure (probability-distance measure). Then,
    \begin{equation*}
        \mathds{E}_{P,Q\sim \mathcal{H}_\theta}[D(P,Q)]
    \end{equation*}
    is non-increasing as $\theta \rightarrow \delta$. In the expectation, $P$ and $Q$ are sampled from $\mathcal{H}_\theta$. 
\end{theorem}

\begin{proof}
%     First, we apply a Lebesgue decomposition to $\lambda$: 
%     \begin{equation*}
%         \lambda = \lambda_{AC} + \lambda_S + \lambda_{D}
%     \end{equation*}
%     where $\lambda_{AC} << \mu$ and $\mu$ is the Lebesgue measure, $\lambda_S \perp \mu$, and $\lambda_D$ is a discrete measure. 

%     The theorem follows by lemma 1, propositions 2,3,4, and the fact that the sum of finitely many decreasing functions is decreasing. (?)

%     \textcolor{red}{INCOMPLETE}
 \end{proof}
% 


%--------------------------------------------------------------
%     EXTRA STUFF BELOW
% --------------------------------------------------------------


\pagebreak

  \begin{definition}[Rashomon Set $\hat{R}_{set}(\mathcal{F}, \theta)$]
        Given 
        \begin{itemize}
            \item $n$ data points $S = \{z_i\}_{i=1}^n$ with $z_i=(x_i,y_i)\sim \ i.i.d \ \mathcal{D}$ on bounded set $\mathcal{Z} = \mathcal{X}\times \mathcal{Y} \subset \mathbb{R}^p \times \mathbb{R}$
            \item Hypothesis space $\mathcal{F} = \{f:\mathcal{X}\rightarrow \mathcal{Y}\}$
            \item (Uniform) prior $\rho$ over $\mathcal{F}$
            \item Loss $\phi:\mathcal{Y}\times \mathcal{Y}\rightarrow \mathbb{R}^+$
            \item Empirical risk $\hat{L}(f) = \frac{1}{n}\sum_i \phi(f(x_i),y_i)$
        \end{itemize}

        A \emph{Rashomon Set} is then defined as $\hat{R}_{set}(\mathcal{F}, \theta) = \{f \in \mathcal{F} \mid \hat{L}(f) \leq \hat{L}(\hat{f})+\theta\}$; $\hat{f}$ is an empirical risk minimizer on $S$ w.r.t $\phi$. 
    \end{definition}

\bigskip


\section{Extra Theorems}

 \begin{theorem}{(Easy Combinatorial Version)}
     Given a binary classification dataset $\mathcal{D}=\mathcal{X}\times \mathcal{Y}$ with $D$ features and $N$ instances, setting $\mathcal{H}$ as the hypothesis class of all $2^N$ binary decision functions, and $\theta < .5$, the $\theta$-Rashomon sets $\mathcal{R}(\mathcal{D};\mathcal{H},\theta)$ computed with respect to classification error are such that 
     \begin{equation*}
         \frac{\mid \mathcal{R}(\mathcal{D};\mathcal{H},\theta) \mid}{\mid \mathcal{H} \mid} 
     \end{equation*}
     is non-increasing as $\theta$ is decreased.
 \end{theorem}

\begin{proof}
    A value of $\theta$ corresponds to an error rate. The number of models in $\mathcal{H}$ with a given error rate is finite since the hypothesis space is finite. Supposing that the error rate $\theta$ is the proportion $m/n$, the number of models with error rate $\theta$ can be represented by the expression ${n \choose m}$ since there are $n \choose m$ ways to make $m$ errors on a dataset with $n$ instances. Thus the proportion of models with error rate $\theta$ is simply ${n \choose m}/{n}$. Reduction of the error rate $\theta$ corresponds to a reduction of $m$. Therefore the theorem will follow by showing that ${n \choose m}/n$ is non-increasing as as $m$ is decreased (and $n$ is fixed).
    \begin{gather*}
        \frac{{n\choose m}}{n} = \frac{\frac{n!}{m!(n-m)!}}{n} = \frac{n!}{nm!(n-m)!} = \frac{(n-1)!}{m!(n-m)!}
    \end{gather*}
    Since $n>m$ and $m \rightarrow 0$ the dominant term in the expression is $(n-m)!$. Since $n$ is fixed, this increases as $m\rightarrow 0$. It follows that $\frac{{n\choose m}}{n}$ decreases as $m$ is decreased. 
\end{proof}

\pagebreak



\begin{proposition}{1}
     Suppose that $\delta=\theta$, and $\mathcal{H}_\theta$ is defined as was previously. Assume that for every $\theta$ the ground truth model $\mathcal{G}$ lies in $\mathcal{H}_\theta$. Let $\lambda$ be the Dirac measure centered on some $h^* \neq \mathcal{G}$, so that $\lambda(E)=1$ when $h^*\in E$ and $\lambda(E)=0$ otherwise. Then:
    \begin{equation*}
        \mathds{E}_{P \sim \mathcal{H}_\theta}[\text{KL}(P || \mathcal{G})]
    \end{equation*}
    is non-increasing as $\theta \rightarrow 0$. In the expectation, $P$ is  sampled from $\mathcal{H}_\theta$. 
\end{proposition}

\begin{proof}
    \begin{equation*}
        \mathds{E}_{P \sim \mathcal{H}_\theta}[\text{KL}(P || \mathcal{G})] = \int_{H_\theta} \text{KL}(P||\mathcal{G}) d\lambda(P) \\
        = \int_{H_\theta} \left[ \int_\mathcal{X} \log\left(\frac{P(x)}{\mathcal{G}(x)}\right) dP(x)\right] d\lambda(P)
    \end{equation*}
    Assuming that the inner integral is $\lambda$-measurable, by definition of Dirac measure the above reduces to $\text{KL}(h^* || \mathcal{G}) \cdot \delta_{h^* \in \mathcal{H}_\theta}$. Now note the following important relation between KL Divergence and cross entropy ($\text{CrsEnt}$):
    \begin{gather*}
        \text{KL}(h^* || \mathcal{G}) = \text{CrsEnt}(h^*,\mathcal{G}) - \text{CrsEnt}(h^*,h^*) \\
        = L(h^*) - \mathds{E}[-\log h^*(\mathcal{X})]
    \end{gather*}
    It follows from the previous two steps that
    \begin{equation*}
        \mathds{E}_{P \sim \mathcal{H}_\theta}[\text{KL}(P || \mathcal{G})] = (L(h^*) - \mathds{E}[-\log h^*(\mathcal{X})])\cdot \lambda(\mathcal{H}_\theta)
    \end{equation*}
    By Lemma 1, this is non-increasing as $\theta$ is decreased.
    
\end{proof}


\begin{proposition}{2}
    The conclusion of proposition 1 holds in the more general case that $\lambda$ is a discrete probability measure   on $\mathcal{H}$ supported on a finite subset. 
\end{proposition}

\begin{proof}
    Suppose $\lambda(h) = p_1 \delta_{h=h_1}(h) + p_2 \delta_{h=h_2}(h) + \ldots + p_K \delta_{h=h_k}(h)$ with $h_1,\ldots,h_K$ distinct hypotheses in $\mathcal{H}$ and $\sum_k p_k=1$. Then the expected KL Divergence over $\mathcal{H}_\theta$ is 
    \begin{equation*}
        \sum_{k=1}^K p_k \text{KL}(h_k || \mathcal{G}) \delta_{h=h_k}(\mathcal{H}_\theta)
    \end{equation*}
    By Lemma 1, each summand is decreasing as $\theta \rightarrow 0$. The sum of finitely many decreasing functions is itself a decreasing function, so the above is decreasing as $\theta \rightarrow 0$. 
\end{proof}

\begin{proposition}{3}
    The conclusion of proposition 1 holds in the case that $\lambda$ is a.c. w.r.t. Lebesgue ($\lambda << \mu$). 
\end{proposition}

\begin{proof}
    Since $\lambda << \mu$, we may apply the Radon-Nikodym theorem. Let $f:\Sigma(\mathcal{H}) \rightarrow [0,\infty)$ be the Radon-Nikodym derivative (here $\Sigma(\mathcal{H})$ is the sigma-algebra generated by $\mathcal{H}$). Then by the Radon-Nikodym theorem, we have
    \begin{gather*}
        \mathds{E}_{P \sim \mathcal{H}_\theta}[\text{KL}(P || \mathcal{G})] = \int_{H_\theta} \text{KL}(P||\mathcal{G}) d\lambda(P) = \int_{H_\theta} \text{KL}(P||\mathcal{G}) f(P) d\mu(P) \\ 
        = \int_{H_\theta} (L(P) - Ent(P)) f(P)d\mu
    \end{gather*}
    \textcolor{red}{INCOMPLETE}
    % BWOC suppose there exists $\theta'< \theta$ such that $\mathds{E}_{P' \sim \mathcal{H}_{\theta'}}[\text{KL}(P' || \mathcal{G})] > \mathds{E}_{P \sim \mathcal{H}_\theta}[\text{KL}(P || \mathcal{G})]$. It follows that:
    % \begin{gather*}
    %     \int_{H_{\theta}} \text{KL}(P||\mathcal{G}) d\lambda(P) - \int_{H_{\theta'}} \text{KL}(P||\mathcal{G}) d\lambda(P) < 0 \\
    % \end{gather*}
    % From the argument in Lemma 1 we know that $H_{\theta'} \subseteq H_{\theta}$, and so by the additivity of integration over unions of disjoint sets $\int_{H_{\theta}} \text{KL}(P||\mathcal{G}) d\lambda(P) = \int_{H_{\theta'}} \text{KL}(P||\mathcal{G}) d\lambda(P) + \int_{H_{\theta}\setminus H_{\theta'}} \text{KL}(P||\mathcal{G}) d\lambda(P)$. It follows that:
    % \begin{equation*}
    %     \int_{H_{\theta}\setminus H_{\theta'}} \text{KL}(P||\mathcal{G}) d\lambda(P) < 0
    % \end{equation*}
    % This is a contradiction because KL-Divergence is non-negative. Hence it must be the case that for $\theta' < \theta$ we have $\mathds{E}_{P' \sim \mathcal{H}_{\theta'}}[\text{KL}(P' || \mathcal{G})] \leq \mathds{E}_{P \sim \mathcal{H}_\theta}[\text{KL}(P || \mathcal{G})]$
\end{proof}

\begin{proposition}{4}
    The conclusion of proposition 1 holds in the case that $\lambda$ is singular w.r.t. Lebesgue $\lambda \perp \mu$. 
\end{proposition}

\begin{proof}
    \textcolor{red}{INCOMPLETE}
\end{proof}


% \begin{proposition}{4}
%     Let $\lambda$ be \emph{any} probability measure on $\mathcal{H}$ and let $\delta,\theta$, $\mathcal{H}_\theta$ be as defined previously, and $D$ be a divergence measure (probability distance metric). Then,
%     \begin{equation*}
%         \mathds{E}_{P,Q\sim \mathcal{H}_\theta}[D(P,Q)]
%     \end{equation*}
%     is non-increasing as $\theta \rightarrow \delta$. In the expectation, $P$ and $Q$ are sampled from $\mathcal{H}_\theta$. 
% \end{proposition}

% \begin{proof}
%     It suffices to show that for any $\theta' < \theta$, we have 
%     \begin{equation*}
%         \mathds{E}_{P,Q\sim \mathcal{H}_{\theta'}}[D(P,Q)] - \mathds{E}_{P,Q\sim \mathcal{H}_\theta}[D(P,Q)] \leq 0
%     \end{equation*}
%     Since $H_\theta \supseteq H_{\theta'}$ we may instead check that
%     \begin{equation*}
%         \mathds{E}_{P,Q\sim \mathcal{H}_\theta\setminus \mathcal{H}_{\theta'}}[D(P,Q)] \geq 0
%     \end{equation*}
%     This follows immediately from the non-negativity of metrics. 
% \end{proof}

\end{document}